\section{Discussion and outlook}
\label{sec:discussion}

In this work we have given an exact construction of the Hamiltonian of mean force
for a system linearly coupled to a Gaussian environment in the commuting (QND)
sector. The central methodological step was to reorganise the reduced equilibrium
problem so that the two sources of difficulty—bath statistics and system operator
structure—are exposed separately. Starting from the Euclidean influence
functional, the bath was eliminated exactly, yielding a bilocal imaginary-time
self-interaction. A Hubbard--Stratonovich transformation then converted this
nonlocal structure into a family of quenched, $\tau$-dependent imaginary-time
propagators acting solely on the system Hilbert space. In the commuting case,
this construction collapses to a closed-form operator expression for
$H_{\mathrm{MF}}(\beta)$.

Several points are worth emphasising. First, the obstruction to writing down a compact Hamiltonian of mean force is
not, in itself, strong coupling or low temperature. Within the present
formulation, both enter only parametrically through the bath kernel (or
equivalently the spectral density). Once the influence functional has been
written, the strength of the coupling can be rescaled against the auxiliary
field introduced by the HS transformation without altering the structural form
of the result. What matters instead is whether the imaginary-time generator
induced by the bath can be represented within a closed operator algebra.

Second, the quenched representation provides a transparent way of diagnosing
this obstruction. For fixed realisations of the auxiliary field, the reduced
equilibrium object is generated by a $\tau$-dependent Hamiltonian on the system
alone. The failure of this generator to collapse to a local, time-independent
operator is therefore an algebraic issue: it arises from the noncommutativity of
the system Hamiltonian with the coupling operator, which forces genuinely
time-ordered structures to survive in the logarithm defining
$H_{\mathrm{MF}}(\beta)$. In the commuting sector, by contrast, time ordering is
trivial and the mean-force Hamiltonian follows directly.

This observation clarifies the relationship between the quantum and classical
theories. In the classical setting the Hamiltonian is a function on phase space,
and all observables commute. The classical Hamiltonian of mean force is therefore
always well defined and local, and the present result reduces directly to the
classical expression obtained by integrating out the bath degrees of freedom.
From a more structural perspective, the same conclusion follows from the Koopman
representation of classical dynamics, in which evolution acts on a commutative
algebra of multiplication operators. The absence of operator noncommutativity
ensures closure. The obstruction encountered here is thus uniquely quantum: it
is the non-Abelian operator algebra generated by $H_Q$ and $f$, rather than the
bath or the temperature, that prevents a local Hamiltonian of mean force from
existing in general.

Finally, the methodology developed here is not restricted to the commuting case.
The quenched representation provides a basis for more systematic analysis of noncommuting couplings, reaction-coordinate
extensions, and controlled enlargements of the operator algebra. In particular,
one may ask whether augmenting the system with additional degrees of freedom
leads to closure at finite order, and how this relates to the construction of
effective Lindblad or mean-force generators in extended spaces. These questions,
together with explicit noncommuting examples, will be taken up in a sequel.

In this sense, the present paper should be read less as an endpoint than as a
reframing. By separating bath spectral data from system algebraic structure, it
makes precise where the difficulty of the Hamiltonian of mean force really lies,
and provides a concrete framework within which that difficulty can be addressed.
